\section{The quadratic-field package}
\label{sec:qp-package}

This appendix supplies the number-field and analytic machinery used in the
main algorithm.  Its input fields always have degree two or four.  An
important distinction is that the enumeration algorithms are polynomial in
the \emph{numerical} discriminant and norm cutoff, not polynomial in their
bitlengths.  This distinction permits exact, potentially very large,
fundamental units without hiding their storage cost.

We use three classical theorems in their following precise forms.  First,
Minkowski's ideal-class bound gives an integral representative of every
ideal class of a degree-$d$ field of discriminant $\Delta$ with norm at most
$(4/\pi)^{r_2}d!d^{-d}\sqrt{|\Delta|}$.  In degrees two and four this is at
most $\sqrt{|\Delta|}$ \cite{Neukirch1999}.  Second, exact LLL reduction of
a fixed-dimensional integer lattice has polynomial bit complexity and,
for $\delta=3/4$, gives
\begin{equation}
 |\mu_{ij}|\leq\tfrac12,
 \qquad \|b_i^*\|\geq2^{-(i-1)/2}\|b_1\|.
 \label{eq:qp-lll}
\end{equation}
The basis transformation is integral and unimodular
\cite{LLL1982}.  Third, if $\chi_p$ is the quadratic character of an odd
prime field and $f$ is not a constant times a square, then
\begin{equation}
 \left|\sum_{z\in\mathbb F_p}\chi_p(f(z))\right|
 \leq (s-1)\sqrt p,
 \label{eq:qp-weil}
\end{equation}
where $s$ counts the distinct roots of $f$ in a splitting field
\cite{Schmidt1976}.  The splitting-field qualification matters.  Every
application below is to a product of at most sixteen linear factors over
$\mathbb F_p$, with an odd root multiplicity.  These classical inputs,
unique factorization of nonzero ideals in rings of integers, and elementary
exact integer arithmetic are the external mathematical prerequisites.
The varying-field Burgess estimate needed here is proved below; it is not
inferred from a fixed-field estimate with an unspecified discriminant
constant.

\subsection{A precise enumeration interface}

For squarefree $d\ne1$ write $F=\mathbb Q(\sqrt d)$ and
$D=|\Delta_F|$.  A supplied root $s^2=d\pmod N$, for odd $N$ coprime to
$d$, gives a ring homomorphism
$\phi_{F,N}:\mathcal O_F\longrightarrow\mathbb Z/N\mathbb Z$:
substitute $s$ for $\sqrt d$ and invert 2 when necessary.  For a
biquadratic field with radicands $d_1,d_2,d_3$, compatible roots mean
\[
 s_i^2=d_i,\qquad s_1s_2=\eta_{12}s_3\pmod N,
 \quad d_3=d_1d_2/g^2,\quad g=\gcd(|d_1|,|d_2|),
\]
where $\eta_{12}=g$ except that $\eta_{12}=-g$ when both $d_1,d_2$ are
negative.  This is the sign of the product of the canonical real or
positive-imaginary radicals.  Substitution, with denominators dividing
four, then gives $\phi_{E,N}$.

\begin{theorem}[Enumeration in degree two and four]
\label{thm:qp-enumeration}
Let $K$ be quadratic or biquadratic, let $D_K=|\Delta_K|$, and supply
the compatible modular roots defining $\phi_{K,N}$.  Let $Z\geq1$ and
assume $N>\max(64D_K,Z)$ is odd and coprime to the defining radicands.
There is a deterministic algorithm of bit cost
\[
 \operatorname{poly}(D_K,Z,\log N)
\]
which either returns a proper divisor of $N$, or returns a finite list of
units modulo $N$ containing the images of a generating set of
$\mathcal O_K^*$ and, for every principal integral ideal $\mathfrak B$ of
norm at most $Z$ coprime to $N$, the image of at least one generator of
$\mathfrak B$.  The number of listed images has the same polynomial bound.
\end{theorem}

The algorithm does not compute a class group.  The following proof also
explains why a large real fundamental unit is compatible with the stated
complexity.

\subsubsection{Real quadratic reduction with exact coefficients}

For a real quadratic discriminant $D$, every nonzero integral ideal is
uniquely $mI(a,b)$, where $m\geq1$,
\begin{equation}
 I(a,b)=a\mathbb Z+\frac{b+\sqrt D}{2}\mathbb Z,
 \quad b\equiv D\pmod2,\quad 4a\mid D-b^2,
 \quad N(I(a,b))=a.
 \label{eq:qp-primitive}
\end{equation}
Here $b$ may be stored modulo $2a$.  To see the assertion, remove the
largest integer scalar from the ideal and write its two-dimensional
Hermite normal form in the basis $1,(D+\sqrt D)/2$.  Closure under
multiplication by the second basis element forces the second diagonal
entry to divide both entries in the first row.  Primitivity makes that
diagonal entry one, giving \eqref{eq:qp-primitive}.  Its remaining closure
condition is exactly the displayed divisibility.

Put $f_D=\lfloor\sqrt D\rfloor$.  A pair is reduced when
\begin{equation}
 0<\sqrt D-b<2a<\sqrt D+b.
 \label{eq:qp-reduced}
\end{equation}
It then has $1\leq a,b\leq f_D$.  For each $b$ the interval for $a$ has
length $b$, so there are at most $f_D(f_D+1)/2<D$ reduced pairs.  Define
\begin{equation}
 a'=(D-b^2)/(4a),\qquad
 b'=f_D-((f_D+b)\bmod 2a').
 \label{eq:qp-gauss}
\end{equation}
The inequalities in \eqref{eq:qp-reduced} give
$\sqrt D-b<2a'<\sqrt D+b$, with the corresponding normalized $b'$;
more explicitly, $2a'\leq f_D+b$ ensures the quotient in
$f_D+b=2a'v+r$ is positive, and $0\leq r<2a'$ gives
$0<\sqrt D-b'<2a'<\sqrt D+b'$.  The inverse is recovered from
$b=f_D-((f_D+b')\bmod2a')$ and $a=(D-b^2)/(4a')$.
Thus \eqref{eq:qp-gauss} permutes the reduced pairs.  Its ideal identity is
\begin{equation}
 \theta I(a,b)=a I(a',b'),
 \qquad \theta=(-b+\sqrt D)/2,
 \qquad |N\theta|=aa'.
 \label{eq:qp-step-identity}
\end{equation}

Start at $a_0=1$, with $b_0$ the largest integer at most $f_D$ of the
parity of $D$.  This is the unique reduced pair for $\mathcal O_F$.
Its cycle has length $L<D$.  If
$\eta_0=1$ and $\eta_{k+1}=\eta_k\theta_k/a_k$, then
$I(a_k,b_k)=\eta_k\mathcal O_F$ and
$0<\sigma_1(\theta_k/a_k)<1$.  Hence
$\varepsilon=\eta_L^{-1}$ is a positive unit greater than one.

We verify that it is fundamental and that this cycle contains every
reduced principal ideal.  Suppose a reduced principal ideal has a
generator $\gamma$ with $0<\sigma_1(\gamma)<1$.  Let
$(a_-,b_-)$ be its predecessor and
$\gamma_-=(a_-/\theta_-)\gamma$.  Clearly
$\sigma_1(\gamma_-) >\sigma_1(\gamma)$.  In fact it is at most one.
Otherwise
\[
 \sigma_1(\gamma)>
 \frac{\sqrt D-b_-}{2a_-}=
 \frac{2a}{\sqrt D+b_-},
 \qquad
 |\sigma_2(\gamma)|<\frac{\sqrt D+b_-}{2}.
\]
Writing
$\gamma_-=a_-v+w(b_-+\sqrt D)/2$ in the predecessor ideal gives
$\gamma=wa+v(-b_-+\sqrt D)/2$.  The two embedding bounds imply
$|v|\sqrt D<1+(\sqrt D+b_-)/2<2\sqrt D$, so
$v\in\{-1,0,1\}$.  If $v=0$, $\sigma_1(\gamma)$ is an integer.
If $v=1$, the inequality $0<\sigma_1(\gamma)<1$ implies $w\leq0$,
so $\sigma_2(\gamma)\leq-(\sqrt D+b_-)/2$.  If $v=-1$, it implies
$w\geq1$, so $\sigma_2(\gamma)>(\sqrt D+b_-)/2$.  All three
possibilities contradict the preceding bounds.

The inverse multipliers range over a finite set of numbers greater than
one.  Repeated inverse steps must therefore reach
$\sigma_1(\gamma_-)=1$, at which point the algebraic integer itself is
one.  This proves that the initial ideal is on the principal cycle.
It also proves that every unit with first embedding in $(0,1)$ is a
negative power of $\varepsilon$.  A generator of any reduced principal
ideal can first be multiplied by a power of the already constructed
$\varepsilon$ and by a sign to enter that interval.  Consequently
\begin{equation}
 \mathcal O_F^*=\{\pm\varepsilon^j:j\in\mathbb Z\},
 \qquad \frac25<\log\varepsilon<\frac12D\log D.
 \label{eq:qp-unit-bounds}
\end{equation}
For the upper bound, each inverse multiplier is
$(\sqrt D+b_k)/(2a_{k+1})<\sqrt D$ and there are fewer than $D$.
For the lower bound write $\varepsilon=(x+y\sqrt D)/2$ with
$|\sigma_2(\varepsilon)|=\varepsilon^{-1}<1$: then $x,y\geq1$ and
$\varepsilon\geq(1+\sqrt5)/2>e^{2/5}$.

To reduce an arbitrary primitive ideal with $a\leq Z$, while $a>f_D$
choose $\widetilde b\equiv b\pmod{2a}$ with
$|\widetilde b|\leq a$ and apply \eqref{eq:qp-step-identity} with
$a'=|D-\widetilde b^2|/(4a)$.  Since $D<a^2$, this is at most $a/4$.
After $O(\log Z)$ steps normalize
$b=f_D-((f_D-b)\bmod2a)$.  Now
$0<\sqrt D-b<2a$.  If the other reduced inequality fails, then
$2a>\sqrt D+|b|$; one further step gives
$2a'=(D-b^2)/(2a)<\sqrt D-|b|\leq\sqrt D$, and its normalized
successor is reduced.  Comparing that pair with the principal cycle
decides principality.  Exact products of the known multipliers recover a
generator when it exists.

The number of factors in these exact products is $O(D+\log Z)$, and
their numerator and denominator bitlengths are $O(\log(D+Z))$.
Quadratic-field rational arithmetic therefore has polynomial bit cost in
$D$ and $\log(Z+2)$.  Multiplying a recovered generator by a suitable
power of $\varepsilon$ gives a balanced generator satisfying
\begin{equation}
 \left|\log\left|\frac{\sigma_1(\gamma)}{\sigma_2(\gamma)}\right|\right|
 <3\log\varepsilon.
 \label{eq:qp-balanced-quadratic}
\end{equation}
The required exponent and resulting coefficient bitlength remain
polynomial in $D,\log(Z+2)$.  Logs can be approximated with certified
error and a fixed slack: after rounding the balancing exponent, a
bounded error gives the factor three in
\eqref{eq:qp-balanced-quadratic}.  Cancellation in a small embedding is
avoided by computing its absolute value as the absolute rational norm
divided by the larger embedding.

Enumerate $a\leq Z$, $b\pmod{2a}$ satisfying
\eqref{eq:qp-primitive}, and $m$ with $m^2a\leq Z$.  This proves the
real quadratic enumeration assertion.  For an imaginary quadratic field,
enumerate integral elements $(x+y\sqrt{-D})/2$ with
$x^2+Dy^2\leq4Z$.  There are polynomially many pairs; they include all
units and all generators of ideals with norm at most $Z$.

\subsubsection{Biquadratic enumeration by parity and weighted lattices}

Write $E=\mathbb Q(\sqrt{d_1},\sqrt{d_2})$, let $d_3$ be the
squarefree product, and put
$L=\mathbb Z+\sum_{i=1}^3\mathbb Z\sqrt{d_i}$.
Then $4\mathcal O_E\subset L$.  Indeed, for
$x=a_0+\sum_i a_i\sqrt{d_i}\in\mathcal O_E$, its rational trace
gives $4a_0\in\mathbb Z$, while each relative trace to the corresponding
quadratic subfield gives $4a_i\in\mathbb Z$.
Test the 256 cosets of $L$ in $\frac14L$ for integrality by the
characteristic polynomial of multiplication.  A monic integral
characteristic polynomial is both necessary and sufficient here.
Hermite normal form gives an integral basis and exact multiplication
tables.  Since the discriminant of $L$ is $256d_1d_2d_3$ and
$[\mathcal O_E:L]\leq256$, we have
\[
 |d_i|\leq256D_E,\qquad |\Delta_{F_i}|\leq1024D_E,
 \qquad D_E=|\Delta_E|.
\]
An integral basis obtained in this way has embeddings of absolute value
at most $49\sqrt{D_E}$.  These deliberately coarse bounds suffice.

Enumerate the Hermite normal forms of all sublattices of index at most
$Z$ and test closure under the integral-basis multiplication tables.
These are exactly the integral ideals of norm at most $Z$.
There are $Z^{O(1)}$ candidates in dimension four.  Fix one,
$\mathfrak B$, of norm $B$.

For each real quadratic subfield $F_i$ compute
\[
 \mathfrak n_i=N_{E/F_i}(\mathfrak B)
   =(\mathfrak B\tau_i(\mathfrak B))\cap\mathcal O_{F_i},
 \qquad N_{F_i/\mathbb Q}(\mathfrak n_i)=B,
\]
where $\tau_i$ fixes $F_i$.  The identity follows prime ideal by prime
ideal in the split, inert, and ramified cases.  Multiplication and
contraction are fixed-dimensional lattice operations.  If any
$\mathfrak n_i$ is nonprincipal, $\mathfrak B$ is nonprincipal.
Otherwise compute a balanced generator $\gamma_i$ and a fundamental
unit $\varepsilon_i$ by the quadratic algorithm.  There are three real
subfields in the totally real case, and one in the CM case.

If $\mathfrak B=(\beta)$, write
\[
 N_{E/F_i}(\beta)=\pm\gamma_i\varepsilon_i^{2k_i+e_i},
 \qquad e_i\in\{0,1\},
 \qquad \beta'=\beta\prod_i\varepsilon_i^{-k_i}.
\]
The relative norm of $\varepsilon_i$ to its own subfield is its square;
to a different quadratic subfield it is its rational norm $\pm1$.
Consequently
$|N_{E/F_i}(\beta')|=|\gamma_i\varepsilon_i^{e_i}|$.
Only the at most eight parity tuples remain unknown.

Let $s_{ik}=\pm1$ indicate which real embedding of $F_i$ is induced by
the $k$th embedding of $E$, choosing the positive embedding of
$\varepsilon_i$ first.  Put
\[
 \ell_i=\log\varepsilon_i,\quad
 \delta_i=\log|\sigma_{i,1}(\gamma_i)/\sigma_{i,2}(\gamma_i)|,
 \quad
 T_k=-\frac14\sum_i s_{ik}(\delta_i+2e_i\ell_i).
\]
For the true parity tuple,
\begin{equation}
 \log|\sigma_k(\beta')|=\tfrac14\log B-T_k,
 \qquad |T_k|=O(D_E\log(4D_E)).
 \label{eq:qp-target}
\end{equation}
In the totally real case use
$\prod_iN_{E/F_i}(\beta')=\beta'^2N_{E/\mathbb Q}(\beta')$.
In the CM case use
$\beta'\overline{\beta'}=N_{E/F^+}(\beta')$.
Both identities give \eqref{eq:qp-target}; conjugate complex embeddings
have equal weights, and the sum of all four $T_k$ is zero.

Weight the Minkowski embedding by $e^{T_k}$, using $\sqrt2$ times the
real and imaginary coordinates for one embedding in each complex pair.
The vector $\beta'$ has length $2B^{1/4}$.  Every nonzero
$x\in\mathfrak B$ has $|Nx|\geq B$, because $(x)\subset\mathfrak B$.
Arithmetic--geometric mean and the zero total weight give the same lower
bound $2B^{1/4}$ for its weighted length.  Thus $\beta'$ is a shortest
vector.

Approximate the weighted basis by a rational basis $Q$ with
$\|(Q-V)V^{-1}\|\leq1/1000$, clear denominators, and apply exact LLL.
In this approximating metric a true shortest vector has length at most
$1.01\|b_1\|$.  Inequalities \eqref{eq:qp-lll}, followed by backward
substitution in Gram--Schmidt coordinates, give successively
\[
 |c_4|\leq2,\qquad |c_3|\leq3,\qquad |c_2|\leq3,
 \qquad |c_1|\leq5.
\]
It is enough to search the fixed box $[-8,8]^4$.  Convert its points to
exact elements of $\mathfrak B$ and retain those of absolute norm $B$.
They generate $\mathfrak B$; a principal ideal has at least one because
its true parity tuple was tried.  For $\mathfrak B=\mathcal O_E$, retain
\emph{all} norm-one candidates and all real-subfield fundamental units.
Every unit is a product of such a balanced candidate and subfield units,
by the parity normalization.  This includes the roots of unity and
proves the full unit-generation assertion.

For completeness, the approximation uses polynomially many bits.
The unweighted embedding matrix of an ideal HNF has entries polynomially
bounded in $D_E,Z$ and determinant of absolute value $B\sqrt{D_E}$.
The adjugate formula bounds its inverse polynomially as well.  After
weighting, the logarithms of both matrix norms are
$O(D_E\log(4D_E)+\log(Z+2))$.  Entrywise interval precision of this
order, with a sufficiently large absolute constant, certifies the
required relative operator error.  Square roots and logarithms are
computed by exact interval arithmetic.  For exponentials, scale the
argument to a bounded interval, use its convergent series, and square
with an explicit propagated error bound.  These operations have
polynomial bit complexity at this precision.  Exact LLL and the
subsequent exact norm computations therefore have the asserted
polynomial cost.  No unit-cost real-arithmetic operation is assumed.

Finally, if $1<\gcd(B,N)<N$, return this factor; since $B\leq Z<N$,
a gcd larger than one cannot equal $N$.  A generator with norm coprime
to $N$ has unit image, by its multiplication-matrix adjugate or the norm
identity.  Integral-basis denominators divide four and are invertible
modulo odd $N$.  This proves Theorem~\ref{thm:qp-enumeration}.

\subsection{Bounded tests remove the ideal-class cocycle}
\label{subsec:qp-characters}

Let $q$ be an odd prime or the product of two distinct odd primes, and
write $\chi$ for the Legendre or Jacobi character modulo $q$.
For each field $K$, let $\mathfrak p_{K,j}$ be the kernel of its
embedding to $\mathbb F_{p_j}$ and put
$\mathfrak q_K=\prod_{p_j\mid q}\mathfrak p_{K,j}$.
This is a product of selected degree-one primes, not necessarily the
whole ideal $q\mathcal O_K$.

\begin{lemma}[Field-specific character assembly]
\label{lem:qp-assembly}
For every required quadratic field $F$ and biquadratic triangle field
$E$, put $M_K=\lceil\sqrt{|\Delta_K|}\rceil$.  Fix $Y\geq2$.  Assume
that all these bounds and the following cutoffs are smaller than $P^-(q)$:
\[
 Z_F=\max(M_F^3,M_FY),\qquad
 Z_E=\max\bigl(M_E^3,M_E\max_{F\subset E}M_F^2\bigr).
\]
Suppose $\chi(a)=1$ for $1\leq a\leq Y$, and that $\chi\circ\phi_K$
is one on all units and all nonzero algebraic integers of absolute norm
at most the corresponding cutoff.  Then it extends to completely
multiplicative functions $\chi_K^*$ on integral ideals, zero exactly on
ideals divisible by selected primes over $q$, with these properties:
\begin{enumerate}
\item $\chi_K^*((\gamma))=\chi(\phi_K(\gamma))$;
\item every prime ideal in a quadratic field above a rational prime
      $r\leq Y$ has value one;
\item for $F\subset E$ and ideals $\mathfrak c$ coprime to the rational
      integer $q$, one has
      $\chi_E^*(\mathfrak c\mathcal O_E)=\chi_F^*(\mathfrak c)$.
\end{enumerate}
\end{lemma}

\begin{proof}
Choose in each inverse ideal class $c^{-1}$ an integral representative
$\mathfrak a_{K,c}$ of norm at most $\sqrt{|\Delta_K|}$, with
$\mathfrak a_{K,1}=\mathcal O_K$.  It is coprime to $q$.
For an ideal $\mathfrak c$ avoiding $\mathfrak q_K$, define
\[
 \chi_K^*(\mathfrak c)=\chi(\phi_K(\gamma)),
 \qquad (\gamma)=\mathfrak a_{K,[\mathfrak c]}\mathfrak c,
\]
and set it to zero otherwise.  Unit triviality makes it independent of
the generator.  For two ideal classes $c_1,c_2$, put
$\mathfrak e=\mathfrak a_{K,(c_1c_2)^{-1}}$.  The identity
\[
 (\mathfrak a_{c_1}\mathfrak c_1)
 (\mathfrak a_{c_2}\mathfrak c_2)(\mathfrak a_{c_1c_2}\mathfrak e)
 =
 (\mathfrak a_{c_1c_2}\mathfrak c_1\mathfrak c_2)
 (\mathfrak a_{c_1}\mathfrak a_{c_2}\mathfrak e)
\]
is an identity between products of principal ideals.  Its two correcting
principal factors have norms at most $M_K^2$ and $M_K^3$, so their
characters are one.  Cancellation proves multiplicativity; the zero
extension is multiplicative because selected prime divisibility is.

A degree-one prime above $r\leq Y$ has its correcting principal ideal
of norm at most $M_Fr$, hence has value one.  This includes ramified
primes.  An inert prime ideal is $(r)$ and has value $\chi(r)=1$.
For compatibility put $\mathfrak b=\mathfrak a_{F,[\mathfrak c]}$.
Extending $(\gamma)=\mathfrak b\mathfrak c$ to $E$ gives
\[
 \chi_E^*(\mathfrak b\mathcal O_E)
 \chi_E^*(\mathfrak c\mathcal O_E)=\chi_F^*(\mathfrak c).
\]
The principal correction defining the first factor has norm at most
$M_E N_F(\mathfrak b)^2\leq M_E M_F^2$, so that factor is one.
\end{proof}

Only finite lists are needed to enforce these hypotheses.  A generating
set of units implies triviality on every unit, including inverse powers.
Any two generators of one principal ideal differ by a unit.  Thus the
finite output of Theorem~\ref{thm:qp-enumeration} enforces every assertion
in Lemma~\ref{lem:qp-assembly}; no class representatives or ideal
characters are computed by the algorithm.

\subsection{Sparse Gram positivity, with ramification retained}
\label{subsec:qp-gram}

Let $S$ be a finite set of distinct rational squareclasses, viewed in a
finite $\mathbb F_2$-space, and put $H=|S|\geq2$.
For every nonzero difference $a+b$, supply its quadratic field
$F_{a+b}$ and the character from Lemma~\ref{lem:qp-assembly}.
For every triangle of distinct vertices supply the compatible
biquadratic field of its three edge classes.  Assume every prime ramified
in an edge field is at most $Y$, and that all the small-prime characters
in that lemma are one.  For $(n,q)=1$ put
\[
 b_\ell(n)=\sum_{N_F\mathfrak a=n}\chi_F^*(\mathfrak a),
 \qquad F=F_\ell.
\]
Each $b_\ell$ is multiplicative.

\begin{theorem}[Ramification-tolerant sparse Gram inequality]
\label{thm:qp-gram}
There is a multiplicative artificial diagonal $b_0$, with
$0\leq b_0(n)\leq\tau(n)$, such that the matrix
$B(n)=(b_{a+b}(n))_{a,b\in S}$ is positive semidefinite for every
$(n,q)=1$.  For $Y$-smooth such $n$, $B(n)-J$ is positive
semidefinite, where $J$ is the all-ones matrix.  Consequently, writing
$\Psi_q(X,Y)$ for the number of $Y$-smooth integers at most $X$ coprime
to $q$,
\begin{equation}
 \max_{\ell\in(S+S)\setminus\{0\}}
 \sum_{\substack{n\leq X\\(n,q)=1}}b_\ell(n)
 \geq \Psi_q(X,Y)-\frac{X(1+\log X)}{H-1}.
 \label{eq:qp-gram-average}
\end{equation}
\end{theorem}

\begin{proof}
First let $r>Y$ be a prime not dividing $q$.  No edge field ramifies.
Choose $h_a=\pm1$ for its local squareclass at $r$, so $F_{a+b}$
splits precisely when $h_a=h_b$.  In a split field, the two prime-ideal
signs are $\alpha$ and $\chi(r)\alpha$.  If $\chi(r)=1$, they are
equal.  For three vertices of the same signature, the four primes of
the compatible biquadratic field show that
\[
 \alpha_{a+b}\alpha_{b+c}\alpha_{a+c}=\chi(r)=1.
\]
Indeed the three extended quadratic primes have pairs of biquadratic
prime signs whose product is the product of all four signs times a
square; the former is $\chi_E^*((r))=\chi(r)$.
Choosing an anchor in each signature class now gives
$\alpha_{a+b}=v_av_b$ for signs $v_a$ within that class.  Only these
triangles, not a full character family, are used.

For $e\geq1$ define the diagonal
\[
 b_0(r^e)=
 \begin{cases}
 e+1,&\chi(r)=1,\\
 0,&\chi(r)=-1,\ e\text{ odd},\\
 1,&\chi(r)=-1,\ e\text{ even}.
 \end{cases}
\]
The split Euler coefficient is
$\sum_{j=0}^e\alpha^j(\chi(r)\alpha)^{e-j}$; the inert coefficient
is zero for odd $e$, and $\chi(r)^{e/2}$ for even $e$.
Thus the four possible matrices are
\[
\begin{array}{c|c}
\text{case}&B(r^e)\\\hline
\chi(r)=1,\ e\text{ odd}&
 (e+1)\bigoplus_Cv_Cv_C^{\mathsf T}\\
\chi(r)=1,\ e\text{ even}&
 \frac{e+2}{2}J+\frac e2hh^{\mathsf T}\\
\chi(r)=-1,\ e\text{ odd}&0\\
\chi(r)=-1,\ e=2m&J\ (m\text{ even}),\quad hh^{\mathsf T}\ (m\text{ odd}).
\end{array}
\]
All are positive semidefinite.

Now let $r\leq Y$, allowing ramification and $r=2$.
The local squareclass group has order four at odd $r$ and order eight
at 2.  Its two-element subgroup $U$ consists of the trivial and the
unramified quadratic class; at 2 these are represented by 1 and 5.
Partition the vertices according to cosets of $U$, and let $v_0,v_1$
be the indicator vectors of the two classes within one coset.
The quotient of two equal local classes is split, of the two different
classes in one coset is inert, and of classes in different cosets is
ramified.  Their coefficients are respectively $e+1$, zero or one
according to the parity of $e$, and one.  All prime-ideal signs are one
by Lemma~\ref{lem:qp-assembly}.  The cross-coset entries of
$B(r^e)-J$ vanish.  Within a coset its contribution is
\[
 \frac{e-1}{2}(v_0+v_1)(v_0+v_1)^{\mathsf T}
 +\frac{e+1}{2}(v_0-v_1)(v_0-v_1)^{\mathsf T}
 \quad(e\text{ odd}),
\]
or $e(v_0v_0^{\mathsf T}+v_1v_1^{\mathsf T})$ for even $e$.
These are positive semidefinite, including repeated local classes.

Multiplicativity makes $B(n)$ the entrywise product of its prime-power
matrices; $B(1)=J$.  The Schur product theorem preserves positive
semidefiniteness.  It also preserves the stronger assertion for smooth
integers, since
$(J+C)\circ(J+D)-J=C+D+C\circ D$ is positive semidefinite when
$C,D$ are.  Therefore
\[
 \sum_{a,b\in S}b_{a+b}(n)
 \geq H^2\mathbf1_{n\text{ is }Y\text{-smooth}}.
\]
Subtract the diagonal and sum; use
$\sum_{n\leq X}\tau(n)\leq X(1+\log X)$.
There are $H(H-1)$ ordered off-diagonal entries.  Dividing by this
number proves \eqref{eq:qp-gram-average}, even when differences repeat.
\end{proof}

\subsection{A uniform quadratic-field Burgess bound}
\label{subsec:qp-burgess}

\begin{theorem}[Uniform prime-or-semiprime orbit estimate]
\label{thm:qp-burgess}
Let $F=\mathbb Q(\sqrt d)$ have discriminant absolute value $D$.
Let $q$ be an odd prime or a product $p_1p_2$ of distinct odd primes,
and supply a root $s^2\equiv d\pmod q$.  Let
$\phi_{F,q}:\mathcal O_F\longrightarrow\mathbb Z/q\mathbb Z$ be
the evaluation map defined by $\sqrt d\mapsto s$ as above; its reduction
modulo each prime dividing $q$ is a degree-one evaluation map.  In the
applications, this root is supplied by the compatible-root construction.
Write $\chi=(\cdot/q)$ for the quadratic character modulo $q$, and suppose
\[
 P^-(q)>D,\qquad D\leq q^{1/200},\qquad
 \chi(\phi_{F,q}(u))=1\quad(u\in\mathcal O_F^*).
\]
For an integral ideal $\mathfrak b$ of norm $B_0\leq\sqrt D$ and
$q^{3/10}\leq W\leq q^{2/5}$, an effective absolute constant $C_B$
satisfies
\begin{equation}
 \left|
 \sum_{\substack{\gamma\in(\mathfrak b\setminus\{0\})/\mathcal O_F^*\\
          1\leq|N\gamma|\leq B_0W}}
       \chi(\phi_{F,q}(\gamma))
 \right|
 \leq C_BD^2Wq^{-1/4000}.
 \label{eq:qp-burgess}
\end{equation}
\end{theorem}

We give the full estimate, including the case of a small prime factor
of $q$.  It follows the uniform sixteenth-moment argument in
\cite{CogDet2026}, with the precise splitting-field form of
\eqref{eq:qp-weil}.

\subsubsection{Planar geometry and representative regions}

If a lattice $\Lambda\subset\mathbb R^2$ has covolume $V$ and
$\|w\|_\infty\geq\mu$ for $w\ne0$, put $\rho=2V/\mu$.
For a rectangle of side lengths $S_1,S_2$,
\begin{equation}
 \frac{(S_1-2\rho)(S_2-2\rho)}V
 \leq\#(\Lambda\cap Q)
 \leq\frac{(S_1+2\rho)(S_2+2\rho)}V,
 \label{eq:qp-rectangle}
\end{equation}
where the lower bound requires $S_i>2\rho$.
Also
\begin{equation}
 \#(\Lambda\cap(z+[-S,S]^2))\leq(2S/\mu+1)^2.
 \label{eq:qp-packing}
\end{equation}
Here is a proof independent of lattice shape.  Choose a shortest
Euclidean vector $w_1$, complete it by a vector $w_2$ of least positive
perpendicular height, and subtract a multiple of $w_1$ to make its
projection at most $\|w_1\|/2$.  It is a basis.  With
$h=V/\|w_1\|$, shortest-vector minimality gives
$\|w_2\|^2\leq\|w_1\|^2/4+h^2\leq\|w_2\|^2/4+h^2$,
and hence both basis lengths are less than $2V/\mu$.
The centered fundamental parallelogram lies inside a sup-norm ball of
radius $\rho$.  Its translates cover the shrunken rectangle and fit
inside the enlarged one, proving \eqref{eq:qp-rectangle}.
Disjoint open sup-norm balls of radius $\mu/2$ prove
\eqref{eq:qp-packing}.

Put $R_*=(B_0W)^{1/2}$.  In the imaginary case use the embedding
$\iota(x)=(\Re x,\Im x)$ and the punctured disk of radius $R_*$.
Divide its sum by $|\mathcal O_F^*|\in\{2,4,6\}$ to obtain the orbit
sum.  In the real case each orbit has a unique representative with
$\sigma_1(\gamma)>0$ and
$0\leq\log|\sigma_1(\gamma)/\sigma_2(\gamma)|<2\log\varepsilon$.
Split this last interval into $K_\varepsilon\leq D^2$ intervals
$[k,k+\delta_k)$ of width at most one, using
\eqref{eq:qp-unit-bounds}.  On a strip use
$\iota_k(x)=(e^{-k/2}\sigma_1(x),e^{k/2}\sigma_2(x))$ and the region
\[
 \Omega_k=\{(v_1,v_2):v_1>0, e^{-\delta_k}v_1<|v_2|\leq v_1,
                                   v_1|v_2|\leq R_*^2\}.
\]
Every region lies in $(-2R_*,2R_*)^2$.  For an ideal of norm $C_0$,
the corresponding lattice has
\begin{equation}
 \tfrac12\sqrt D C_0\leq V(C_0)\leq\sqrt D C_0,
 \quad \mu(C_0)\geq\sqrt{C_0/2},
 \quad \rho(C_0)\leq2\sqrt D\sqrt{C_0}.
 \label{eq:qp-ideal-lattice}
\end{equation}
The minimum bounds follow from $|Nx|\geq C_0$ for a nonzero ideal
element, and are unchanged by the real rescaling.

\subsubsection{The grid and bilinear shifts}

Assume first $q\geq2^{8000}$.  Set
\[
 \epsilon_0=\epsilon_1=q^{-1/4000},\quad
 M_g=\lceil2/\epsilon_0\rceil,\quad
 \ell=2R_*/M_g,\quad B=q^{1/32},\quad
 A=\frac{\epsilon_1\ell}{4\sqrt{B_0}B}.
\]
Thus $\frac23\epsilon_0R_*<\ell\leq\epsilon_0R_*$,
$2\rho(\mathfrak b)/\ell<\epsilon_1$, and
\begin{equation}
 \tfrac18q^{0.11825}\leq A\leq\tfrac14q^{0.16825}.
 \label{eq:qp-A-range}
\end{equation}
Tile $[-2R_*,2R_*)^2$ by squares of side $\ell$.
Each contains at most $8\epsilon_0^2W/\sqrt D$ lattice points.
The disk boundary is covered by four graphs of Lipschitz constant one.
The real-strip boundary is covered by four rays, two hyperbola arcs
with slope of absolute value at most one, and the axis segment.
A graph meets at most $6M_g$ square closures; altogether at most
$40M_g$ boundary squares occur.  Their total contribution is at most
\begin{equation}
 960Wq^{-1/4000}
 \label{eq:qp-boundary}
\end{equation}
for each strip.

For an interior square $Q\subset\Omega_k$, let
\[
 \mathcal A=\{\beta\in\mathcal O_F\setminus\{0\}:
       \|\iota(\beta)\|_\infty\leq A,
       \,\phi_{F,q}(\beta)\in(\mathbb Z/q\mathbb Z)^*\},
 \qquad
 \mathcal B_k=\{\omega\in\mathfrak b:
       \|\iota_k(\omega)\|_\infty\leq\sqrt{B_0}B\}.
\]
Here $\iota$ is the unscaled embedding.  The following effective bounds
are uniform:
\begin{equation}
 \frac{A^2}{\sqrt D}<|\mathcal A|<\frac{16A^2}{\sqrt D},
 \qquad
 \frac{2B^2}{\sqrt D}<|\mathcal B_k|<\frac{16B^2}{\sqrt D}.
 \label{eq:qp-shift-counts}
\end{equation}
To check the possible loss of units in the first set, let
$V_0=V(\mathcal O_F)$ and let $\mathcal A_0$ include zero and all
points of its box.  Since $\rho(\mathcal O_F)/A<1/8$,
\eqref{eq:qp-rectangle} gives
$|\mathcal A_0|>49A^2/(16V_0)$ and
$|\mathcal A_0|<8A^2/V_0$.
Every nonzero element has norm below $2A^2<\sqrt q$, so it cannot
belong to the selected prime above $p_2\geq\sqrt q$ (or above a prime
modulus $q$).  For the smaller prime $p_1>D\geq3$, if $p_1\leq A$,
the same rectangle estimate bounds its kernel points by
$81A^2/(80V_0)$; here $p_1\geq5$ and
$\rho(\mathfrak p_1)/A\leq2\sqrt D/\sqrt A<1/8$.
If $p_1>A$, packing bounds them by $16A<A^2/V_0$.
Subtracting proves the first pair of bounds.
The second pair follows from
$\rho(\mathfrak b)/(\sqrt{B_0}B)\leq2\sqrt D/B<1/4$.

For every pair of shifts,
$\|\iota_k(\beta\omega)\|_\infty\leq2A\sqrt{B_0}B
=\epsilon_1\ell/2=:h$.  The unmatched points when translating the
square by $\beta\omega$ lie between its $h$-enlargement and
$h$-shrinking.  Formula \eqref{eq:qp-rectangle} bounds their number by
$16\epsilon_1\ell^2/(\sqrt D B_0)$.
Define
\[
 \nu(z)=\#\{(\gamma,\beta)\in
   (\mathfrak b\cap\iota_k^{-1}(Q))\times\mathcal A:
             \phi(\gamma)\phi(\beta)^{-1}=z\},
 \quad
 T(z)=\sum_{\omega\in\mathcal B_k}\chi(z+\phi(\omega)).
\]
Averaging the translated sums and using
$|\chi(\phi(\beta))|=1$ gives
\begin{equation}
 |S_Q|\leq\frac{\sum_{z\bmod q}\nu(z)|T(z)|}
                  {|\mathcal A||\mathcal B_k|}
             +\frac{16\epsilon_1\ell^2}{\sqrt D B_0}.
 \label{eq:qp-shift-average}
\end{equation}

\subsubsection{Ratio collisions lift to exact identities}

The first moment obeys
\begin{equation}
 \sum_z\nu(z)\leq K_0:=\frac{128A^2\ell^2}{DB_0}.
 \label{eq:qp-first-moment}
\end{equation}
A collision of two ratios implies
$\delta=\gamma_1\beta_2-\gamma_2\beta_1\in\mathfrak b\mathfrak q_F$.
Indeed $\mathfrak b$ and $\mathfrak q_F$ are coprime, because
$B_0\leq\sqrt D<P^-(q)$.  This ideal has norm $B_0q$, whereas
\[
 |N\delta|\leq16B_0WA^2\leq B_0q^{0.7365}<B_0q.
\]
Hence $\delta=0$.  For a fixed pair $(\gamma_1,\beta_2)$, the nonzero
product $\mu=\gamma_1\beta_2$ has absolute norm below $q^2$.
Every matching $\beta_1$ generates an ideal divisor of $(\mu)$.
A quadratic field has at most $\tau(|N\mu|)^2$ such ideal divisors,
since at most two primes lie over a rational prime.

An effective constant $C_*$ satisfies $\tau(n)\leq C_*n^{1/800}$.
For example take the product, over primes below $2^{800}$, of the
maximum of $(j+1)\ell^{-j/800}$ over integers $j\geq0$; each maximum
occurs in a finite effectively bounded range.  Primes at least
$2^{800}$ need no factor because $2^j\geq j+1$.
The number of admissible generators of one ideal divisor is at most
$6\log q$.  In the real case every embedding absolute value of a
generator in the $A$-box lies in $[A^{-1},A]$; two such generators
differ by $\pm\varepsilon^n$ with
$|n|<5\log A<\log q$, using \eqref{eq:qp-unit-bounds}.
The imaginary case has at most six units.
Since $\log q<1000q^{1/1000}$, we obtain
\begin{equation}
 \sum_z\nu(z)^2\leq C_1q^{3/500}K_0,
 \qquad C_1=6000C_*^2.
 \label{eq:qp-second-moment}
\end{equation}

\subsubsection{The sixteenth moment with a small semiprime factor}

For each $p_j\mid q$ let
$\psi_j:\mathcal B_k\longrightarrow\mathbb F_{p_j}$ be the
evaluation map.  If two images agree, their difference lies in
$\mathfrak b\mathfrak p_j$ and has norm at most $8B_0B^2$.
Thus $\psi_j$ is injective whenever $p_j>8B^2$.  This always holds
for the larger factor $p_2$ and for a prime modulus, since
$8B^2=8q^{1/16}<\sqrt q$.  Otherwise each fiber has size at most
\begin{equation}
 \left(\frac{2\sqrt2 B}{\sqrt{p_j}}+1\right)^2
 \leq\frac{32B^2}{p_j},
 \label{eq:qp-fibers}
\end{equation}
by \eqref{eq:qp-packing}.

Expand $\sum_z|T(z)|^{16}$ over ordered sixteen-tuples from
$\mathcal B_k$.  If a residue modulo a prime has odd multiplicity,
\eqref{eq:qp-weil} bounds its complete character sum by $16\sqrt p$;
otherwise use $p$.  The tuple classes are:
\begin{enumerate}
\item Every individual $\omega$ occurs an even number of times.
Pairing the sixteen positions gives at most
$16!|\mathcal B_k|^8<16!16^8q^{1/2}$ tuples, each contributing at most
$q$.
\item Some residue modulo $p_1$ has odd multiplicity.
For a semiprime, some individual element then has odd multiplicity,
and injectivity at $p_2$ gives an odd residue there too.  The complete
sum is at most $256\sqrt q$ across fewer than $16^{16}q$ tuples.
For a prime modulus the same upper bound is valid.
\item Individual multiplicities are not all even, but all residues
modulo $p_1$ have even multiplicity.  Necessarily
$q=p_1p_2$ and $p_1\leq8B^2$.  Pairing and \eqref{eq:qp-fibers}
give at most $16!512^8q p_1^{-8}$ tuples.  Their sums are at most
$16p_1\sqrt{p_2}=16\sqrt q\sqrt{p_1}$, hence total
$O(q^{3/2}p_1^{-15/2})$.
\end{enumerate}
Chinese remaindering multiplies the two prime sums, including zero
character values.  Therefore
\begin{equation}
 \sum_{z\bmod q}|T(z)|^{16}\leq C_2q^{3/2},
 \quad C_2=16!16^8+256\cdot16^{16}+16!16\cdot512^8.
 \label{eq:qp-sixteenth}
\end{equation}
The third case is why no balance assumption on the two prime factors
is needed.

\subsubsection{The power saving and effective constants}

Apply H\"older to
$\nu^{7/8}(\nu^2)^{1/16}(|T|^{16})^{1/16}$, using
\eqref{eq:qp-first-moment}--\eqref{eq:qp-sixteenth} and
\eqref{eq:qp-shift-counts}.  The first term of
\eqref{eq:qp-shift-average} is at most
\[
 64(C_1C_2)^{1/16}\frac{\ell^2}{B_0}
 D^{1/16}
 \left(\frac{B_0}{A^2\ell^2}\right)^{1/16}
 B^{-2}q^{3/32+3/8000}.
\]
Substitute $A=\epsilon_1\ell/(4\sqrt{B_0}B)$ and
$B_0/\ell^2<4\epsilon_0^{-2}q^{-3/10}$.
With $D\leq q^{1/200}$ the resulting exponent of $q$ is
\[
 \frac1{3200}+\frac1{16000}-\frac3{80}
 +\frac1{32000}-\frac{15}{256}+\frac3{32}+\frac3{8000}
 =-\frac1{640}<-\frac1{4000}.
\]
Thus each interior square contributes at most
$(C_3+16)\ell^2B_0^{-1}q^{-1/4000}$, where
$C_3=128(C_1C_2)^{1/16}$.  There are at most
$16B_0W/\ell^2$ squares.  Add \eqref{eq:qp-boundary} and at most
$D^2$ real strips to obtain \eqref{eq:qp-burgess}, for example with
$C_B=\max(32004,16C_3+1216)$.

For $q<2^{8000}$, each orbit gives a distinct integral ideal
$\gamma\mathfrak b^{-1}$ of norm at most $W$.  Hence its count is
at most $\sum_{n\leq W}\tau(n)\leq W(1+\log W)$, which is
absorbed by $32004Wq^{-1/4000}$.  This proves
Theorem~\ref{thm:qp-burgess} with an effective absolute constant.

\subsection{From orbit sums to rationally coprime ideal sums}

\begin{corollary}
\label{cor:qp-ideal-sum}
Assume Lemma~\ref{lem:qp-assembly} and the hypotheses of
Theorem~\ref{thm:qp-burgess}.  For $q^{3/10}\leq W\leq q^{2/5}$,
\[
 U_F(W):=\sum_{N\mathfrak a\leq W}\chi_F^*(\mathfrak a)
 \quad\text{satisfies}\quad
 |U_F(W)|\leq C_BD^3Wq^{-1/4000}.
\]
At $X=q^{1/3}$,
\begin{equation}
 \left|\sum_{\substack{n\leq X\\(n,q)=1}}b_F(n)\right|
 \leq 2C_BD^3Xq^{-1/4000}
       +Xq^{-1/50}(1+\log X).
 \label{eq:qp-coprime-sum}
\end{equation}
\end{corollary}

\begin{proof}
There are at most $D$ ideal classes: choose representatives of norm at
most $\sqrt D$, use at most $\tau(n)\leq n$ ideals of norm $n$, and
sum $n\leq\sqrt D$.  Multiplication by the chosen inverse-class
representative identifies each class sum with an orbit sum in an ideal
of norm at most $\sqrt D$.  The first assertion follows.

A prime modulus exceeds $X$, so its rational divisibility removes
nothing.  If $q=p_1p_2$, the larger factor also exceeds $X$.
At $p_1$, terms divisible by the selected ideal already have character
zero.  Terms divisible by the other degree-one prime ideal sum to its
sign times $U_F(X/p_1)$, by multiplicativity.  If
$p_1\leq q^{1/50}$, then
$q^{47/150}\leq X/p_1\leq q^{1/3}$, within the orbit estimate's
range.  Otherwise the elementary divisor bound gives at most
$Xq^{-1/50}(1+\log X)$.  This proves \eqref{eq:qp-coprime-sum}.
\end{proof}

\subsection{An elementary smooth-number lower bound}
\label{subsec:qp-smooth}

We only need the following weak, effective estimate; no smooth-number
asymptotic or heuristic is used:
\begin{equation}
 \pi(Y)-\pi(Y/2)>\frac{Y}{12\log Y}\qquad(Y\geq4096).
 \label{eq:qp-elementary-primes}
\end{equation}
We include an elementary proof, as in \cite{CogDet2026}; stronger
classical explicit bounds are available in \cite{RosserSchoenfeld1962}.
First $\prod_{r\leq x}r\leq4^x$.  Induct on the integer part of $x$.
An even endpoint greater than two is composite.  For an odd endpoint
$2k+1$, every prime between $k+1$ and $2k+1$ divides
$\binom{2k+1}{k}\leq4^k$; multiply this by the inductive bound
$4^{k+1}$ for the remaining primes.

Now put $n=\lfloor Y/2\rfloor\geq2048$.  In the prime factorization
of $\binom{2n}{n}$, the contribution of each prime power is at most
$2n$.  Primes greater than $\sqrt{2n}$ occur to exponent at most one,
and primes in $(2n/3,n]$ do not occur.  Hence
\[
 \frac{4^n}{2n+1}\leq\binom{2n}{n}
 \leq (2n)^{\sqrt{2n}}4^{2n/3}
                    (2n)^{\pi(2n)-\pi(n)}.
\]
After taking base-two logarithms and dividing by $n$, the two error
terms $\log_2(2n+1)/n$ and
$\sqrt{2n}\log_2(2n)/n$ decrease for $n\geq2048$.
At 2048 their sum is less than $781/2048$, so
\[
 \frac{\pi(2n)-\pi(n)}n\log_2(2n)
 >\frac23-\frac{781}{2048}>\frac14.
\]
Since $n>2047Y/4096$, $\pi(Y/2)=\pi(n)$, and $\log2>0.69$,
this implies \eqref{eq:qp-elementary-primes}.
In particular $\pi(Y)>Y/(12\log Y)$ on this entire range.

If
$Y=\lfloor X^{1/u}\rfloor$ is sufficiently large and
$u\leq Y/(24\log Y)$, distinct products of $u$ primes in that
interval are distinct integers at most $X$.  If $P^-(q)>Y$, none is
divisible by $q$.  Writing $P_Y$ for the number of available primes,
\[
 \Psi_q(X,Y)\geq\binom{P_Y}{u}
 \geq(P_Y/u)^u
 \geq X(25\log X)^{-u}.
\]
For the final inequality use
$Y\geq(4096/4097)X^{1/u}$ once $X^{1/u}\geq4097$, and
$u\log Y\leq\log X$; the constants 12 and 25 leave ample slack.
All thresholds are effective.

\subsection{A conservative explicit scale certificate}

\begin{lemma}[One common numerical threshold]
\label{lem:qp-explicit-scale}
Let $m\geq4096$, $L=2^m$,
$R=2^{\lceil\sqrt{256Lm}\rceil}$, $V=\log R$,
$L\leq\theta<2L$, and $X=e^{\theta/3}$.
Let $u=\lceil(\log X)/V\rceil$ and
$Y_1=\lfloor X^{1/u}\rfloor$.  Then the prime-product hypotheses
above hold, and
\[
 \Psi_q(X,R)\geq Xe^{-V/8}\quad(P^-(q)>R),
 \qquad
 \frac{X(1+\log X)}{\pi(R)}\leq Xe^{-3V/4}.
\]
If $q=e^\theta$ is a permitted integer modulus, $D\leq4R^3$, and
the other hypotheses of Corollary~\ref{cor:qp-ideal-sum} hold, its
right-hand side is at most $Xe^{-\theta/8000}$, strictly less than
$\frac12Xe^{-V/8}$.  In particular $D\leq q^{1/200}$ and
$q\geq2^{8000}$ on this range.
\end{lemma}

\begin{proof}
We give coarse bounds so that the natural-logarithm and binary-scale
conventions are explicit.  They are
\[
 11\sqrt{Lm}<V<17\sqrt{Lm},\qquad
 V^2>121Lm,\qquad V\geq704m.
\]
The first two follow from $0.69<\log2<1$; the last follows from
$2^m\geq m^2$ and $\sqrt m\geq64$ for $m\geq4096$.
The ratio $\sqrt{m/2^m}$ decreases for $m\geq2$.
Thus $3V/\theta\leq51\sqrt{m/2^m}<1/2$ and $u\geq2$.
It follows that
$\log Y_1\geq2V/3-\log2\geq V/2$.
Also $u\leq L$ and
$\log(24VL)\leq2m+7<V/2$.
Consequently $Y_1\geq4096$ and
$u\leq Y_1/(24\log Y_1)$.
Finally
\[
 \frac{u\log(25\log X)}V
 \leq\frac{2(m+3)}{363m}+\frac{m+3}{704m}<\frac18,
\]
which proves the smooth lower bound.  By
\eqref{eq:qp-elementary-primes}, the logarithm of the factor multiplying
$e^{-V}$ in the diagonal bound is at most
$\log(12V(1+2L/3))<3m<V/4$.

The constants in the Burgess proof can also be bounded explicitly.
For $j\geq0$, $(j+1)2^{-j/800}\leq2048$; for example write
$j=800k+r$ and use $k+1\leq2^k$.
Hence $C_*\leq2048^{2^{800}}$ in
\eqref{eq:qp-second-moment}.  The displayed formulas for $C_1,C_2,C_3$
and $C_B$ give the convenient bound
$\log(4C_B)<2^{810}$.
Now $3\log D\leq12V\leq204\sqrt{Lm}$.
At $m=4096$,
\[
 2^{810}+204\sqrt{Lm}<2^{2063}
       <2^{4083}<L/8000.
\]
Both ratios to $L$ decrease thereafter.  Thus
$\log(4C_B)+3\log D\leq\theta/8000$,
and the first term in \eqref{eq:qp-coprime-sum} is at most
$\frac12Xe^{-\theta/8000}$.
The other term satisfies the same bound because
\[
 \log\bigl(2(1+\theta/3)\bigr)\leq m+1
 <L/100\leq\theta(1/50-1/8000).
\]
The inequality $D\leq q^{1/200}$ follows separately from
$4V\leq68\sqrt{Lm}<L/200$, first at $m=4096$ and then by the same
monotonic ratio.  The condition $q\geq2^{8000}$ is immediate.
The preceding budget also implies
$\theta/8000>V/8+\log2$, which proves the final strict comparison.
\end{proof}

The threshold is deliberately conservative.  It is a certificate of
effectivity, not a recommendation to run the numerical guards at that
scale.  The algorithms use an elementary terminating fallback for the
finite range below it.  This concludes the common package.
